\section*{Capítulo \ref{ch:b3:ode} --- Ecuaciones diferenciales
ordinarias}

\begin{solution}{exo:b3:ode:1}
(a) Separando variables: $x(t) = \frac1{1 - t}$ en $\intoo{-\infty}1$:
explosión en $T_+ = 1$, con $x(t) \to
+\infty$ --- exactamente el~\cref{thm:b3:ode:maximal}(2). (b)
$x(t) = \tan t$ en $\intoo{-\pi/2}{\pi/2}$: explosión en ambos extremos.
(c) $x(t) = \frac{1}{1 + \eu^{-t}}$, global: la solución permanece en
$\intoo01$, un conjunto acotado, de modo que el grafo no puede abandonar
todo \omterm{def:b3:topology:compact}{compacto} de $\R\times\R$ en
tiempo finito --- $T_\pm
= \pm\infty$. Obsérvese que (a) y (b) no contradicen
el~\cref{cor:b3:ode:globallinear}: $x^2$ y $1 + x^2$ tienen crecimiento
superlineal.
\end{solution}

\begin{solution}{exo:b3:ode:2}
(a) Es la estimación del~\cref{cor:b3:ode:uniqueness} con $t_0
= 0$. Optimalidad: para $x' = Lx$ (globalmente $L$-lipschitziana), dos
soluciones difieren exactamente en $(x_0 - y_0)\eu^{Lt}$. (b) La
estimación dice: la aplicación de
\omterm{ex:b3:ode:planeclassification}{flujo} en tiempo $t$ es
$\eu^{L\abs t}$-lipschitziana en la condición inicial ---
\omterm{def:b3:topology:continuity}{continuidad}, uniformemente para $t$
en \omterm{def:b3:topology:compact}{compactos}; en conjuntos acotados
con $f$ no globalmente lipschitziana, hágase lo mismo en un tubo
\omterm{def:b3:topology:compact}{compacto} alrededor de las trayectorias
con la constante local.
\end{solution}

\begin{solution}{exo:b3:ode:3}
(a) $\abs{\sin(tx)} \leq 1$: acotada, es decir, crecimiento lineal con
$\alpha = 0$, $\beta = 1$: el~\cref{cor:b3:ode:globallinear} en
$U = \R\times\R$.
(b) $\abs{tx/(1 + x^2)} \leq \abs t\cdot\frac12$: de nuevo sublineal (de
hecho, acotada en los \omterm{def:b3:topology:compact}{compactos}
temporales): global. (c) $X = (x, x')$:
$X' = \bigl(\begin{smallmatrix}0 & 1\\ -q(t)
& 0\end{smallmatrix}\bigr)X$: lineal con coeficientes
\omterm{def:b3:topology:continuity}{continuos}: global (el marco
del~\cref{thm:b3:ode:linearstructure}). (d) Global; Grönwall como en
el~\cref{cor:b3:ode:globallinear}:
$\norm{x(t)} \leq \norm{x(t_0)}\,\eu^{M\abs{t - t_0}}$ con
$M = \sup\vertiii{A}$.
\end{solution}

\begin{solution}{exo:b3:ode:4}
(a) $A^2 = -I$ para la primera: $\eu^{tA} = \cos t\,I + \sin
t\,A = \bigl(\begin{smallmatrix}\cos t & -\sin t\\ \sin t &
\cos t\end{smallmatrix}\bigr)$. Bloque de Jordan: $\lambda I$ y $N$
conmutan: $\eu^{tA} =
\eu^{\lambda t}\bigl(\begin{smallmatrix}1 & t\\ 0 &
1\end{smallmatrix}\bigr)$. Tercera: $A^2 = I$: $\eu^{tA} =
\cosh t\,I + \sinh t\,A$.

(b) Sistema $X' = \bigl(\begin{smallmatrix}0&1\\-1&0
\end{smallmatrix}\bigr)X + \bigl(\begin{smallmatrix}0\\
\cos\omega t\end{smallmatrix}\bigr)$; Duhamel con la
\omterm{thm:b3:ode:linearstructure}{resolvente} de rotación da las
soluciones particulares: para $\omega \neq 1$,
$x_p(t) = \frac{\cos(\omega t)}{1 -
\omega^2}$ (verifíquese directamente); para $\omega = 1$, la integral
$\int_0^t\sin(t - s)\cos s\,\dd s = \frac t2\sin t$ produce el
crecimiento \emph{secular} $x_p = \frac t2\sin t$: resonancia --- el
forzamiento bombea energía a la frecuencia natural y la amplitud crece
linealmente.
\end{solution}

\begin{solution}{exo:b3:ode:5}
(a) $w' = x_1x_2'' - x_1''x_2 = x_1(-px_2' - qx_2) - (-px_1'
- qx_1)x_2 = -p\,w$: $w(t) = w(t_0)\eu^{-\int_{t_0}^tp}$,
nunca nulo o idénticamente nulo. Si $w \neq 0$, los vectores
$(x_i, x_i')(t_0)$ son independientes en $\R^2$ y, como el espacio de
soluciones tiene dimensión $2$ (el~\cref{thm:b3:ode:linearstructure}
para el sistema), $(x_1,
x_2)$ es una base.

(b) Con $u = \int\frac{\eu^{-\int p}}{x_1^2}$: $x_2 = x_1u$,
$x_2' = x_1'u + \frac{\eu^{-\int p}}{x_1}$, y
\[
x_2'' + px_2' + qx_2
= u\,(x_1'' + px_1' + qx_1) +
\Bigl(-\,p\frac{\eu^{-\int p}}{x_1} +
p\frac{\eu^{-\int p}}{x_1}\Bigr) = 0
\]
(los términos cruzados se cancelan exactamente; desarróllese con
cuidado). Para
$t^2x'' - 2x = 0$, i.e.\ $x'' - \frac2{t^2}x = 0$ ($p = 0$)
con $x_1 = t^2$: $u = \int t^{-4} = -\frac1{3t^3}$, de modo que $x_2
= -\frac1{3t}$: solución general $at^2 + \frac bt$.
\end{solution}

\begin{solution}{exo:b3:ode:6}
Equilibrios $0, 1$; $f'(x) = 1 - 2x$: $f'(0) = 1 > 0$ (inestable --- las
soluciones próximas se alejan, como muestra la forma explícita),
$f'(1) = -1 < 0$: asintóticamente \omterm{def:b3:ode:stability}{estable}
(el~\cref{thm:b3:ode:linearization} en dimensión $1$). Para
$x(0) \in \intoo01$: allí $f > 0$, de modo que, mientras la solución
permanezca en $\intoo01$, crece; nunca puede alcanzar $0$ ni $1$
(unicidad: son trayectorias), de modo que permanece, está acotada ---
luego es global --- y crece hacia un límite $L
\in \intoc{x(0)}1$. Si $f(L) \neq 0$, entonces $x' \geq c > 0$ cerca del
límite, lo que fuerza $x$ más allá de $L$: así pues, $f(L) = 0$,
$L = 1$; simétricamente, $x \to 0$ en $-\infty$. Explícitamente, $x(t) =
\frac1{1 + C\eu^{-t}}$ lo confirma todo. En general: si una solución
escalar autónoma tuviera $x'(t_0) = 0$, entonces $x(t_0)$ es un
equilibrio y la unicidad hace $x$ constante; en caso contrario,
$f(x(t))$ mantiene un signo fijo (no se anula nunca y
$t \mapsto f(x(t))$ es \omterm{def:b3:topology:continuity}{continua}):
$x$ es estrictamente monótona --- de modo que una solución periódica no
constante es imposible.
\end{solution}

\begin{solution}{exo:b3:ode:7}
(a) $\frac{\dd}{\dd t}H(x,y) = H_xx' + H_yy' = H_xH_y +
H_y(-H_x) = 0$.
(b) Los conjuntos de nivel de $H = \frac{y^2}2 + \frac{x^4}4$ son
\omterm{def:b3:topology:compact}{compactos} ($H$ coerciva), de modo que
las soluciones quedan atrapadas en
\omterm{def:b3:topology:compact}{compactos}: globales y acotadas
(el~\cref{thm:b3:ode:maximal}). Estabilidad de $(0,0)$: $V = H$ es
definida positiva ($H = 0$ solo en el origen) con $\dot V = 0$:
el~\cref{thm:b3:ode:lyapunov}. La linealización $x' = y$, $y' = 0$ tiene
la matriz nilpotente no diagonalizable de valor propio $0$:
el~\cref{thm:b3:ode:linearization} calla (su hipótesis
$\operatorname{Re}\lambda < 0$ falla) y, en efecto, el sistema
linealizado es inestable ($y_0 \ne 0$ deriva) mientras que el no lineal
es \omterm{def:b3:ode:stability}{estable}: linealizar en un equilibrio
no hiperbólico no demuestra nada.
\end{solution}

\begin{solution}{exo:b3:ode:8}
(a) $\dot V = y\,y' + \sin x\cdot x' = y(-\sin x - cy) +
y\sin x = -cy^2 \leq 0$, y $V = \frac{y^2}2 + (1 - \cos x)$ es definida
positiva en $\{\abs x < 2\pi\}$ alrededor del origen:
\omterm{def:b3:ode:stability}{estable} (el~\cref{thm:b3:ode:lyapunov}).
(b) La matriz linealizada en $(0,0)$ es
$\bigl(\begin{smallmatrix}0 & 1\\ -1 &
-c\end{smallmatrix}\bigr)$, de polinomio característico
$\lambda^2 + c\lambda + 1$: raíces $\frac{-c \pm \sqrt{c^2 -
4}}2$ --- ambas reales negativas si $c \geq 2$, complejas de parte real
$-\frac c2 < 0$ si $0 < c < 2$. En todos los casos
$\operatorname{Re}\lambda < 0$:
el~\cref{thm:b3:ode:linearization} da estabilidad asintótica (pese a la
degeneración $\dot V$). (c) En $(\pi, 0)$: $\sin(\pi + u) = -\sin u$,
linealización
$\bigl(\begin{smallmatrix}0&1\\1&-c\end{smallmatrix}\bigr)$,
característico $\lambda^2 + c\lambda - 1$: raíces de signos opuestos
($\lambda_+\lambda_- = -1$). A lo largo del vector propio inestable, el
sistema lineal tiene la solución explícitamente divergente
$\eu^{\lambda_+t}v_+$ con $\lambda_+ > 0$: el péndulo invertido es
inestable para todo amortiguamiento.
\end{solution}

\begin{solution}{exo:b3:ode:9}
Para $\alpha \in \intoo01$: además de $x \equiv 0$, cada
\[
x_c(t) = \begin{cases}0 & t \leq c,\\
\bigl((1-\alpha)(t - c)\bigr)^{1/(1-\alpha)} & t \geq c,
\end{cases}
\]
es $\mathcal C^1$ y resuelve la ecuación (el exponente
$\frac1{1-\alpha} > 1$ hace que la derivada se anule en $c$): un
\omterm{def:b3:topology:continuity}{continuo} de soluciones que pasan por $(0,0)$. Para $\alpha = 1$: $x
\mapsto \abs x$ es globalmente $1$-lipschitziana
($\abs{\abs a - \abs b} \leq \abs{a - b}$), de modo que
el~\cref{thm:b3:ode:cauchylipschitz} se aplica y la única solución que
pasa por $0$ es $x \equiv 0$. La frontera es exactamente la condición de
Lipschitz local en $0$: $\abs x^\alpha$ tiene allí cocientes
incrementales no acotados para $\alpha < 1$.
\end{solution}

\begin{solution}{exo:b3:ode:10}
(a) Supóngase $\norm{x(t_2)} > R$ para algún $t_2 > 0$ con
$\norm{x(0)} \leq R$, y sea $t_1 = \sup\{t \leq t_2 :
\norm{x(t)} \leq R\}$: entonces $\norm{x(t_1)} = R$ y $\norm{x(t)} > R$
en $\intoc{t_1}{t_2}$. En ese intervalo, $g(t) = \norm{x(t)}^2$ cumple
$g'(t) = 2\langle x(t),
F(x(t))\rangle \leq 0$ (la hipótesis se aplica: $\norm{x(t)}
\geq R$), de modo que $g(t_2) \leq g(t_1) = R^2$: contradicción. La bola
es positivamente invariante. (b) Una solución atrapada en la bola
compacta no puede tener $T_+ < \infty$
(el~\cref{thm:b3:ode:maximal}(2)): global hacia el futuro. Sistema
gradiente: $\frac{\dd}{\dd t}G(x(t)) =
\langle\nabla G, -\nabla G\rangle = -\norm{\nabla G(x(t))}^2
\leq 0$: $G$ decrece, de modo que la solución permanece en $\{G \leq
G(x_0)\}$, que es acotado (coercividad: fuera de una bola grande,
$G > G(x_0)$) y cerrado: \omterm{def:b3:topology:compact}{compacto}. La
salida es imposible: toda solución de un sistema gradiente coercivo es
global hacia el futuro, deslizándose cuesta abajo para siempre.
\end{solution}

\begin{solution}{exo:b3:ode:11}
(a) Lema de comparación: sea $w = x - y$ en el intervalo común;
$w(0) \geq 0$ y $w' = x' - y' \geq F(x) - F(y) =
c(t)w$ con $c(t) = \frac{F(x) - F(y)}{x - y}$ acotada en los intervalos
de tiempo \omterm{def:b3:topology:compact}{compactos} ($F$
\omterm{def:b3:ode:lipschitz}{localmente lipschitziana}); entonces
$(w\eu^{-\int c})' \geq 0$, de modo que $w \geq 0$ en todo el intervalo.
Con $F(x) = x^2$: $y(t) = \frac1{1 - t}$ explota en $1$, y $x
\geq y$ mientras $x$ viva; si $T_+ > 1$, entonces $x$ sería finita en
$t = 1$ dominando a $y \to \infty$: absurdo. $T_+ \leq 1$.

(b) La comparación invertida (el mismo lema, con los papeles
intercambiados): en $\intcc01\cap\intco0{T_+}$, $t^2 \leq 1$ da $x' \leq
x^2 + 1$, mientras que $z(t) = \tan(t + \frac\pi4)$ cumple $z' =
z^2 + 1$, $z(0) = 1 = x(0)$: por tanto, $x \leq z$ mientras ambas estén
definidas. Como $z$ es finita en $\intco0{\frac\pi4}$, $x$ no puede
explotar antes de
$\frac\pi4$: $T_+ \geq \frac\pi4$.

(c) Juntas: $\frac\pi4 \leq T_+ \leq 1$ (numéricamente, $T_+
\approx 0.96$). Moraleja: para $x' = F(t, x)$ con $F$ superlineal en
$x$, las soluciones explotan en tiempo finito siempre que
$\int^\infty\frac{\dd s}{F(s)} < \infty$ (la solución de comparación
alcanza el infinito en ese tiempo finito); el crecimiento lineal, donde
la integral diverge, fuerza la existencia global
(el~\cref{exo:b3:ode:3}). Es la misma integral de Osgood que en
el~\cref{pb:b3:complete:1}, gobernando ahora la huida hacia el infinito
en lugar de la huida desde cero.
\end{solution}

\begin{solution}{exo:b3:ode:12}
(a) $W' = uv'' - u''v = u(-q_2v) - (-q_1u)v = (q_1 -
q_2)\,uv$.

(b) Sean $a < b$ ceros consecutivos de $u$; normalícese $u >
0$ en $\intoo ab$ (de modo que $u'(a) > 0$, $u'(b) < 0$ --- no nulas por
unicidad, pues $u(a) = u'(a) = 0$ forzaría $u
\equiv 0$). Supóngase que $v$ no tiene ningún cero en $\intoo ab$;
normalícese $v > 0$ allí (de donde $v(a), v(b) \geq 0$ por
\omterm{def:b3:topology:continuity}{continuidad}). Intégrese (a):
\[
W(b) - W(a) = \int_a^b(q_1 - q_2)\,uv \;\leq\; 0 .
\]
Pero $W(a) = u(a)v'(a) - u'(a)v(a) = -u'(a)v(a) \leq 0$ y
$W(b) = -u'(b)v(b) \geq 0$: luego $W(b) - W(a) \geq 0$. Igualdad en toda
la cadena: $\int(q_1 - q_2)uv = 0$ con $uv > 0$ en el intervalo abierto
obliga a $q_1 = q_2$ allí; y $W(a) =
W(b) = 0$ obliga a $v(a) = v(b) = 0$; entonces $W \equiv 0$ en
$\intcc ab$ (su derivada se anula), es decir, $(v/u)' = -W/u^2 = 0$ en
$\intoo ab$: $v \propto u$.

(c) Tómense $q_1 = m$ y $u = \sin(\sqrt m(t - t_0))$, cuyos ceros
consecutivos distan $\pi/\sqrt m$, y $q_2 = q \geq
m$: por (b), toda solución $v$ de $v'' + qv = 0$ se anula en cada
intervalo abierto de longitud $\pi/\sqrt m$ (en la alternativa
degenerada $q \equiv m$, $v \propto u$ también se anula). Si, en cambio,
$q \leq 0$: aplíquese (b) con $q_1 = q$, $u =
v$ y $q_2 = 0$ con la solución sin ceros $\mathbf 1$ de $v'' = 0$. Si
$v$ tuviera dos ceros consecutivos, (b) forzaría un cero de $\mathbf 1$
entre ellos o bien el caso degenerado $\mathbf 1 \propto v$ --- ambos
absurdos: $v$ se anula a lo sumo una vez. Ensayos: para $q = 1$,
$\sin t$ se anula cada $\pi
= \pi/\sqrt1$, como se predijo; para $q = -1$, $\sinh t$ se anula
exactamente una vez y $\eu^t$ ninguna --- a lo sumo un cero, como se
predijo.
\end{solution}

\begin{solution}{pb:b3:ode:1}
\textbf{1.} $\dot E = yy' + \sin x\cdot x' = -y\sin x +
y\sin x = 0$: la energía es una integral primera. En una
\omterm{thm:b3:ode:maximal}{solución maximal},
$y^2 = 2(E_0 + \cos x) \leq 2(E_0 + 1)$: $y$ está acotada; entonces
$\abs{x(t)} \leq \abs{x(0)} +
t\sup\abs y$ crece a lo sumo linealmente: en cualquier intervalo finito
de tiempo la trayectoria permanece en un
\omterm{def:b3:topology:compact}{compacto} de $\R^2$, de modo que
el~\cref{thm:b3:ode:maximal}(2) obliga a $T_\pm = \pm\infty$.
Equilibrios $(k\pi, 0)$; linealización $\bigl(\begin{smallmatrix}0&1\\
\mp1&0\end{smallmatrix}\bigr)$ con $-\cos(k\pi) = \mp1$: valores propios
$\pm\iu$ para $k$ par (tipo \omterm{ex:b3:groups:actions}{centro},
inconcluyente por sí solo) y $\pm1$ para $k$ impar (punto de silla).

\textbf{2.} $E$ constante a lo largo de las soluciones confina cada
trayectoria a un conjunto de nivel $\{y^2 = 2(E_0 + \cos x)\}$. Para
$E_0 = -1$: solo los puntos $(2k\pi, 0)$. Para $-1 < E_0 <
1$: escribiendo $E_0 = -\cos a$, el conjunto es una unión disjunta de
curvas cerradas $y = \pm\sqrt{2(\cos x - \cos a)}$ sobre $x \in
[2k\pi - a, 2k\pi + a]$, una alrededor de cada equilibrio
\omterm{def:b3:ode:stability}{estable}. Para $E_0 = 1$: las curvas
$y = \pm2\cos\frac x2$ que unen puntos de silla consecutivos --- la
separatriz --- junto con los propios puntos de silla. Para $E_0 > 1$:
dos grafos $y =
\pm\sqrt{2(E_0 + \cos x)}$, definidos para todo $x$, que nunca tocan
$y = 0$.

\textbf{3.} $V = E + 1 = \frac{y^2}2 + (1 - \cos x)$
se anula en $(0,0)$, es positiva en un
\omterm{def:b3:topology:topology}{entorno} punteado ($\abs x < 2\pi$) y
$\dot V = 0 \leq 0$: el~\cref{thm:b3:ode:lyapunov} da la estabilidad. No
asintótica: la solución que pasa por $(a, 0)$ ($0 < a$ pequeño)
permanece en la curva de nivel $E = -\cos a$, cuya distancia al origen
es positiva (la curva corta el eje $x$ solo en $\pm a$):
$\varphi_t(a, 0) \not\to (0,0)$.

\textbf{4.} En la curva de nivel $C_a$: no hay equilibrios ($y = 0$
obligaría a $x = \pm a$ con $\sin(\pm a) \neq 0$ para $0 < a <
\pi$), de modo que la velocidad $\norm{(y, -\sin x)}$ tiene un mínimo
positivo $m$ en el \omterm{def:b3:topology:compact}{compacto} $C_a$.
Sígase la solución desde $(a, 0)$: en el semiplano inferior
$x' = y < 0$, de modo que $x$ decrece de $a$ a $-a$ en tiempo finito
(las integrales de cuarto y de medio periodo convergen: el análisis de
la pregunta 5), alcanzando $(-a, 0)$; por la simetría
$(x, y) \mapsto (x, -y)$, $t
\mapsto -t$ de la ecuación, el semiplano superior se recorre de vuelta
en el mismo tiempo $\frac T2$: la solución vuelve a $(a, 0)$ en el
instante $T$. La unicidad (el~\cref{cor:b3:ode:uniqueness}) propaga
entonces: $x(t + T) =
x(t)$ para todo $t$: periódica.

\textbf{5.} En la rama donde $y > 0$: $y = \sqrt{2(\cos x
- \cos a)}$ y $\dd t = \frac{\dd x}{y}$; integrando $x$ de $-a$ a $a$ se
obtiene el medio periodo, y la simetría $x
\mapsto -x$ vuelve a partir la integral por la mitad:
\[
T(a) = 4\int_0^a\frac{\dd x}{\sqrt{2(\cos x - \cos a)}} .
\]
Convergencia en $x = a^-$: $\cos x - \cos a = \sin(a)(a - x) +
O((a-x)^2)$ con $\sin a > 0$: el integrando se comporta como
$\bigl(2\sin a\,(a - x)\bigr)^{-1/2}$,
\omterm{def:b3:lebesgue:l1}{integrable}.

\textbf{6.} Con $\sin\frac x2 = k\sin\varphi$, $k =
\sin\frac a2$: $\cos x - \cos a = 2(k^2 - \sin^2\frac x2)
= 2k^2\cos^2\varphi$ y $\frac12\cos\frac x2\,\dd x =
k\cos\varphi\,\dd\varphi$, de modo que
\[
T(a) = 4\int_0^{\pi/2}
\frac{\dd\varphi}{\sqrt{1 - k^2\sin^2\varphi}} .
\]
Cuando $a \to 0^+$, $k \to 0$: para $k \leq k_0 < 1$ el integrando está
dominado por $(1 - k_0^2\sin^2\varphi)^{-1/2}$,
\omterm{def:b3:topology:continuity}{continua} en $\intcc0{\pi/2}$: el
teorema de convergencia dominada da $T \to
4\cdot\frac\pi2 = 2\pi$. Desarrollando $(1 - u)^{-1/2} = 1 +
\frac u2 + \frac{3u^2}8 + \cdots$ con $u =
k^2\sin^2\varphi$ (convergencia normal para $k < 1$) y
$\int_0^{\pi/2}\sin^2 = \frac\pi4$:
\[
T(a) = 2\pi\Bigl(1 + \frac{k^2}{4} + O(k^4)\Bigr) :
\]
el isocronismo solo vale en primer orden; el periodo crece con la
amplitud.

\textbf{7.} Cuando $k \uparrow 1$, los integrandos crecen hacia $(1
- \sin^2\varphi)^{-1/2} = \frac1{\cos\varphi}$, cuya integral diverge:
por convergencia monótona, $T(a) \to
+\infty$ cuando $a \to \pi^-$.

\textbf{8.} En la rama $y = 2\cos\frac x2$ ($\abs x <
\pi$): $\frac{\dd x}{2\cos(x/2)} = \dd t$; con $u = \frac
x2$, $\int\frac{\dd u}{\cos u} =
\ln\tan\bigl(\frac u2 + \frac\pi4\bigr)$, de modo que $t =
\ln\tan\bigl(\frac x4 + \frac\pi4\bigr)$, es decir,
\[
x(t) = 4\arctan(\eu^t) - \pi,
\qquad y(t) = x'(t) = \frac{4\eu^t}{1 + \eu^{2t}} =
\frac{2}{\cosh t} .
\]
Verificación con la energía: con $\theta = \arctan\eu^t$,
$\sin2\theta = \frac1{\cosh t}$, de modo que $\cos x = -\cos4\theta =
-1 + \frac{2}{\cosh^2t}$ y $E = \frac{y^2}2 - \cos x =
\frac2{\cosh^2t} + 1 - \frac2{\cosh^2t} = 1$: la trayectoria está sobre
la separatriz, y derivar $y^2 = 2(1 +
\cos x)$ donde $y > 0$ reproduce $y' = -\sin x$. Límites: $x
\to \pm\pi$ y $y \to 0$ cuando $t \to \pm\infty$.

\textbf{9.} La \omterm{def:b3:groups:action}{órbita} tiende al punto de
silla $(\pi, 0)$ hacia el futuro y a $(-\pi, 0)$ hacia el pasado, pero
no llega nunca: si alcanzara $(\pi, 0)$ en un instante finito $t^*$, dos
soluciones maximales distintas --- la
solución separatriz y la solución constante en el punto de silla ---
pasarían por el mismo punto $(t^*, (\pi, 0))$, en contra
del~\cref{cor:b3:ode:uniqueness}. A los puntos de silla solo se llega
asintóticamente.

\textbf{10.} Para $E_0 > 1$: $y^2 = 2(E_0 + \cos x) \geq
2(E_0 - 1) > 0$: $y$ mantiene su signo, y $\abs{x'} = \abs y
\geq \sqrt{2(E_0 - 1)}$: $x$ es estrictamente monótona y global
(pregunta 1), con $x(t) \to \pm\infty$. Como $y(t) =
\pm\sqrt{2(E_0 + \cos x(t))}$ y $\cos$ es $2\pi$-periódica, $y$ recupera
su valor cada vez que $x$ avanza $2\pi$; el tiempo necesario es
\[
\tau(E_0) = \int_{x_0}^{x_0 + 2\pi}\frac{\dd
x}{\sqrt{2(E_0 + \cos x)}} =
\int_{-\pi}^{\pi}\frac{\dd x}{\sqrt{2(E_0 + \cos x)}}
\]
(sustitución; periodicidad): $y$ es $\tau$-periódica --- el péndulo gira
con velocidad de rotación asintóticamente constante
$2\pi/\tau \approx \sqrt{2E_0}$ a energías grandes.

\textbf{11.} El método, por orden: la \emph{energía} ($\dot E =
0$) reduce el \omterm{ex:b3:ode:planeclassification}{flujo}
bidimensional a curvas de nivel unidimensionales; la \emph{acotación} de
$y$ en cada nivel, más la \omterm{thm:b3:ode:maximal}{salida de los compactos}, da la existencia
global; la \emph{\omterm{def:b3:topology:compact}{compacidad}} de los niveles cerrados da cotas de
velocidad y, con ellas, la periodicidad; la \emph{unicidad} convierte el
primer retorno en periodicidad exacta, prohíbe la llegada en tiempo
finito a los puntos de silla y separa los tipos de
\omterm{def:b3:groups:action}{órbita}; la \emph{linealización y
Lyapunov} clasifican los equilibrios; la \emph{integral del periodo} se
analiza con los teoremas de convergencia del~\cref{ch:b3:lebesgue}. En
ningún momento poseímos --- ni necesitamos --- una solución general en
forma cerrada: la teoría cualitativa extrajo de la propia ecuación todos
los rasgos del movimiento.

\textbf{12.} Por partes: $W_n = \int_0^{\pi/2}\sin^{2n-1}\varphi
\cdot\sin\varphi\,\dd\varphi = (2n-1)\int_0^{\pi/2}
\sin^{2n-2}\varphi\cos^2\varphi\,\dd\varphi = (2n-1)(W_{n-1}
- W_n)$, de modo que $W_n = \frac{2n-1}{2n}W_{n-1}$; con $W_0 =
\frac\pi2$ y $\prod_{j=1}^n\frac{2j-1}{2j} =
\frac{(2n)!}{4^n(n!)^2}$ (sepárese $(2n)! = 2^nn!\prod(2j-1)$), la
inducción da el valor exhibido. Serie binómica: $(1 -
u)^{-1/2} = \sum_{n\geq0}c_nu^n$ con $c_n =
\frac{(2n)!}{4^n(n!)^2} \in \intoc01$, radio $1$. Para $k
\leq k_0 < 1$ y $u = k^2\sin^2\varphi$, la serie
$\sum c_nk^{2n}\sin^{2n}\varphi$ converge normalmente en $\varphi$
($c_nk^{2n} \leq k_0^{2n}$), de modo que la integración término a
término en la fórmula de la pregunta 6 es lícita:
\[
T(a) = 4\sum_{n\geq0}c_nk^{2n}W_n
= 4\cdot\frac\pi2\sum_{n\geq0}c_n^2\,k^{2n}
= 2\pi\sum_{n\geq0}
\Bigl(\frac{(2n)!}{4^n(n!)^2}\Bigr)^{\!2}k^{2n} .
\]
Con $c_0 = 1$, $c_1 = \frac12$, $c_2 = \frac38$: $T(a) =
2\pi\bigl(1 + \frac{k^2}4 + \frac{9k^4}{64} + O(k^6)\bigr)$, con resto
uniforme para $k \leq k_0$ (cola dominada por una serie geométrica).

\textbf{13.} $k = \sin\frac a2 = \frac a2 - \frac{a^3}{48} +
O(a^5)$, de modo que
\[
k^2 = \frac{a^2}4 - \frac{a^4}{48} + O(a^6),
\qquad
k^4 = \frac{a^4}{16} + O(a^6) ,
\]
y
\[
\frac{T(a)}{2\pi} = 1 + \frac14\Bigl(\frac{a^2}4 -
\frac{a^4}{48}\Bigr) + \frac9{64}\cdot\frac{a^4}{16} +
O(a^6)
= 1 + \frac{a^2}{16} + \frac{11\,a^4}{3072} + O(a^6) ,
\]
puesto que $-\frac1{192} + \frac9{1024} = \frac{-16 + 27}{3072}
= \frac{11}{3072}$.

\textbf{14.} En la forma elíptica de la pregunta 6, $a \mapsto
k = \sin\frac a2$ es una biyección
\omterm{def:b3:topology:continuity}{continua} estrictamente creciente de
$\intoo0\pi$ sobre $\intoo01$, y para cada $\varphi \in \intoc0{\pi/2}$
el integrando $(1 -
k^2\sin^2\varphi)^{-1/2}$ es estrictamente creciente en $k$: $T$ es
estrictamente creciente.
\omterm{def:b3:topology:continuity}{Continuidad}: en $k \leq k_0 < 1$ el
integrando está dominado por la
\omterm{def:b3:topology:continuity}{continua} $(1 -
k_0^2\sin^2\varphi)^{-1/2}$, de modo que la convergencia dominada se
aplica a lo largo de $k' \to k$. Con los límites $T \to 2\pi$ cuando $a
\to 0^+$ (pregunta 6) y $T \to +\infty$ cuando $a \to \pi^-$ (pregunta
7), la monotonía estricta y el teorema del valor intermedio hacen de $T$
una biyección de $\intoo0\pi$ sobre
$\intoo{2\pi}{+\infty}$.

\textbf{15.} Un reloj cuenta oscilaciones; regulado a amplitud
infinitesimal, contabiliza el periodo armónico $2\pi$ por oscilación (en
la unidad de tiempo del péndulo). Funcionando con amplitud $a$, el
periodo verdadero es $T(a) = 2\pi\bigl(1 + \frac{a^2}{16} +
O(a^4)\bigr)$: el reloj contabiliza $2\pi$ mientras transcurre realmente
$T(a)$, de modo que atrasa la fracción
\[
\frac{T(a) - 2\pi}{T(a)} = \frac{a^2}{16} + O(a^4) .
\]
Para $a = 0.2$ rad (unos $11.5$ grados): $\frac{a^2}{16} =
\frac{0.04}{16} = \frac1{400}$, y un día tiene $86400$ s: el reloj
pierde $86400/400 = 216$ segundos --- unos tres minutos y medio --- al
día. De ahí los dos remedios históricos: imponer una amplitud minúscula
y constante (el escape), o doblar la ligadura para que el periodo sea
exactamente independiente de la amplitud (las mejillas cicloidales de
Huygens, 1657).

\textbf{16.} En $\tau(E_0) = \int_{-\pi}^{\pi}\frac{\dd
x}{\sqrt{2(E_0 + \cos x)}}$, el integrando es, para cada $x$ fijo,
estrictamente decreciente en $E_0$: $\tau$ es estrictamente decreciente.
Cuando $E_0 \downarrow 1$, los integrandos crecen puntualmente hacia
$\bigl(2(1 + \cos x)\bigr)^{-1/2} =
\frac1{2\abs{\cos\frac x2}}$, cuya integral sobre $\intoo{-\pi}\pi$
diverge ($\cos\frac x2$ se anula a primer orden en $\pm\pi$): la
convergencia monótona da $\tau(E_0)
\to +\infty$. Cuando $E_0 \to +\infty$:
\[
\sqrt{2E_0}\,\tau(E_0) = \int_{-\pi}^{\pi}
\frac{\dd x}{\sqrt{1 + \cos x/E_0}} \longrightarrow 2\pi
\]
por convergencia dominada (para $E_0 \geq 2$, el integrando es a lo sumo
$\sqrt2$). Así, $\tau \approx 2\pi/\sqrt{2E_0}$: una vuelta lleva el
tiempo de la rotación libre a velocidad $\sqrt{2E_0}$, reducido el
potencial a un rizo --- en consonancia con la velocidad de rotación de
la pregunta 10.

\textbf{17.} Los ejes soportan las soluciones explícitas $t
\mapsto (x_0\eu^{\alpha t}, 0)$ y $t \mapsto (0,
y_0\eu^{-\gamma t})$, junto con el equilibrio $(0, 0)$: son uniones de
\omterm{def:b3:groups:action}{órbitas}. El campo es $\mathcal
C^1$, luego localmente lipschitziano; una
solución que arrancara en $Q$ y tocara un eje pasaría por un punto de
una de esas \omterm{def:b3:groups:action}{órbitas} y, por
el~\cref{cor:b3:ode:uniqueness}, coincidiría con ella --- imposible,
viviendo una sobre el eje y la otra no. Así pues, $Q$ es invariante en
ambos sentidos del tiempo. Equilibrios en $Q$: $x > 0$ obliga a
$\alpha - \beta y = 0$ y $y > 0$ obliga a $\delta x - \gamma = 0$: el
único punto
$(x_*, y_*) = \bigl(\frac\gamma\delta,
\frac\alpha\beta\bigr)$.

\textbf{18.} A lo largo de una solución,
\[
\dot H = \Bigl(\delta - \frac\gamma x\Bigr)x' +
\Bigl(\beta - \frac\alpha y\Bigr)y'
= (\delta x - \gamma)(\alpha - \beta y) +
(\beta y - \alpha)(\delta x - \gamma) = 0 .
\]
$f(x) = \delta x - \gamma\ln x$ tiene $f'' = \gamma/x^2 > 0$, $f'$ que
solo se anula en $x_*$ y $f \to +\infty$ tanto en $0^+$ como en
$+\infty$: estrictamente convexa y propia, con mínimo $f(x_*)$;
análogamente $g(y) = \beta y - \alpha\ln y$, con mínimo $g(y_*)$. Así,
$H \geq h_*$ con igualdad solo en $(x_*,
y_*)$, y cada subnivel $\{H \leq h\}\cap Q$ es
\omterm{def:b3:topology:compact}{compacto}: $f(x) \leq h - g(y_*)$
confina $x$ a un intervalo \omterm{def:b3:topology:compact}{compacto} de
$\intoo0{+\infty}$ por propiedad, y otro tanto $y$, y el conjunto es
cerrado en $\R^2$ puesto que $H \to +\infty$ en la frontera de $Q$. Una
\omterm{thm:b3:ode:maximal}{solución maximal} permanece en su conjunto
de nivel \omterm{def:b3:topology:compact}{compacto}, de modo que no
puede abandonar todo \omterm{def:b3:topology:compact}{compacto} en
tiempo finito: el~\cref{thm:b3:ode:maximal} la hace global.

\textbf{19.} Fíjese $h > h_*$ y póngase $c = h - g(y_*) >
f(x_*)$. Como $f$ decrece estrictamente de $+\infty$ a $f(x_*)$ en
$\intoc0{x_*}$ y crece estrictamente de vuelta hasta $+\infty$ en
$\intco{x_*}{+\infty}$, la ecuación $f(x) = c$ tiene exactamente dos
raíces $x_- < x_* < x_+$, y $\{f \leq c\} =
\intcc{x_-}{x_+}$. Para $x \in \intoo{x_-}{x_+}$: $g(y) = h -
f(x) > g(y_*)$ tiene exactamente dos raíces $y_-(x) < y_* <
y_+(x)$, \omterm{def:b3:topology:continuity}{continuas} en $x$ (inversas
de las restricciones \omterm{def:b3:topology:continuity}{continuas}
estrictamente monótonas de $g$ a cada lado de $y_*$), con
$y_\pm(x) \to y_*$ cuando $x \to x_\pm$; en $x =
x_\pm$ la única solución es $y = y_*$. Así, $\{H = h\}\cap
Q$ es la unión de los grafos de $y_+$ y $y_-$ sobre $\intcc{x_-}{x_+}$,
pegados en $(x_\pm, y_*)$: una curva cerrada alrededor de $(x_*, y_*)$
--- el análogo de los óvalos del péndulo.

\textbf{20.} Sea $C_h = \{H = h\}\cap Q$ con $h > h_*$: el único
equilibrio de $Q$ queda fuera de $C_h$, de modo que el campo no se anula
nunca sobre él. Signos: $x' = \beta x(y_* - y)$, $y' =
\delta y(x - x_*)$: el movimiento va hacia la derecha por debajo de la
recta $y
= y_*$, hacia arriba a la derecha de $x = x_*$, hacia la izquierda por
encima y hacia abajo a la izquierda --- circulación en sentido
antihorario. Sígase la solución desde un punto $(x_0, y_-(x_0))$ de la
rama inferior abierta, donde $x' > 0$: el tiempo hasta alcanzar la
esquina derecha $B = (x_+, y_*)$ es
\[
\int_{x_0}^{x_+}\frac{\dd x}{\beta x\,\bigl(y_* -
y_-(x)\bigr)} .
\]
Cerca de $x_+$, elíjase $x' \in \intoo{x_*}{x_+}$; para $x \in
\intcc{x'}{x_+}$, $f(x_+) - f(x) \geq f'(x')\,(x_+ - x)$ ($f'$ es
creciente y positiva pasado $x_*$), mientras que la relación de nivel y
la desigualdad de Taylor dan $g(y_-(x)) -
g(y_*) \leq \frac12\,\bigl(\max g''\bigr)\,(y_* -
y_-(x))^2$ en el rango \omterm{def:b3:topology:compact}{compacto} de $y$
de $C_h$: por tanto, $y_* -
y_-(x) \geq c\,\sqrt{x_+ - x}$ con $c > 0$, y el integrando es
$O\bigl((x_+ - x)^{-1/2}\bigr)$: \omterm{def:b3:lebesgue:l1}{integrable}
--- la convergencia de la pregunta 5, trasladada. En el resto de la rama
el integrando es \omterm{def:b3:topology:continuity}{continuo}. Así, $B$
se alcanza en tiempo finito; allí $y' = \delta y_*(x_+ - x_*) > 0$, la
\omterm{def:b3:groups:action}{órbita} entra en la región $x > x_*$,
$y > y_*$, asciende hasta la esquina superior $(x_*, y_+)$ por la
estimación simétrica (intercambiando los papeles de $f$ y $g$), y así
por los cuatro arcos: tras un tiempo finito $T > 0$ la solución vuelve a
su punto de partida. Por el~\cref{cor:b3:ode:uniqueness} es
$T$-periódica --- el argumento de la pregunta 4, palabra por palabra.

\textbf{21.} En una \omterm{def:b3:groups:action}{órbita} $T$-periódica
de $Q$, $t \mapsto \ln
x(t)$ es $\mathcal C^1$ y $T$-periódica, de modo que
\[
0 = \int_0^T(\ln x)'\,\dd t
= \int_0^T(\alpha - \beta y)\,\dd t
= \alpha T - \beta\int_0^Ty\,\dd t ,
\]
lo que da $\frac1T\int_0^Ty = \frac\alpha\beta$; análogamente, $0 =
\int_0^T(\ln y)' = \delta\int_0^Tx\,\dd t - \gamma T$ da
$\frac1T\int_0^Tx = \frac\gamma\delta$. Las medias temporales son los
valores de equilibrio, en toda \omterm{def:b3:groups:action}{órbita} y
con independencia de la amplitud --- una ley de conservación que nadie
introdujo a mano.

\textbf{22.} Con captura, el sistema vuelve a ser de tipo
Lotka--Volterra, con parámetros $\alpha -
\varepsilon$, $\beta$, $\gamma + \varepsilon$, $\delta$ (el equilibrio
\omterm{def:b3:topology:interior}{interior} persiste porque
$\varepsilon <
\alpha$). La pregunta 21 aplicada al nuevo sistema:
\[
\bar x = \frac{\gamma + \varepsilon}{\delta}
\ \ (\text{la media de presas sube}),
\qquad
\bar y = \frac{\alpha - \varepsilon}{\beta}
\ \ (\text{la media de depredadores baja}) :
\]
la captura indiscriminada desplaza el equilibrio a favor de las presas.
Los datos de d'Ancona se leen al revés: la guerra recortó la pesca,
$\varepsilon$ bajó, de modo que la media de depredadores
$(\alpha - \varepsilon)/\beta$ subió y la de presas bajó --- una mayor
proporción de tiburones en las capturas, exactamente lo que registraron
las lonjas del Adriático. Este es el principio de Volterra, el mismo
mecanismo que hay detrás de las paradojas de los plaguicidas: diezmar
ambos niveles tróficos beneficia al nivel que es devorado.

\textbf{23.} La receta, en ambos casos: (i) una integral primera --- $E$
para el péndulo, $H$ aquí, hallada separando $\dd y/\dd x$ --- colapsa
el plano en curvas; (ii) la propiedad de ser propios y la \omterm{def:b3:topology:compact}{compacidad} de
los conjuntos de nivel dan la existencia global por
el~\cref{thm:b3:ode:maximal}; (iii) la geometría de los niveles ---
óvalos, por la forma de $\cos$ allí y la convexidad estricta de $f$ y
$g$ aquí --- se lee en la integral, no en el
\omterm{ex:b3:ode:planeclassification}{flujo}; (iv) un campo que no se
anula en un óvalo \omterm{def:b3:topology:compact}{compacto}, más
singularidades \omterm{def:b3:lebesgue:l1}{integrables} en las esquinas,
fuerza un tiempo de retorno finito; (v) la unicidad
(el~\cref{cor:b3:ode:uniqueness}) convierte el retorno en periodicidad y
prohíbe la llegada en tiempo finito a los equilibrios; (vi) los
dividendos --- desarrollos del periodo, leyes de medias --- provienen de
los teoremas de convergencia aplicados a las integrales resultantes. Ni
$x'' = -\sin x$ ni Lotka--Volterra admiten una solución elemental en
forma cerrada (integrales elípticas en un caso, curvas de nivel
trascendentes en el otro), y en ningún momento hizo falta: la propia
ecuación, interrogada cualitativamente, entregó el movimiento entero.

\textbf{24.} En un nivel de energía $E_0 > 1$, $y^2 = 2(E_0 +
\cos x) > 0$: extremos de $y^2$ en $\cos x = \pm1$, lo que da las cotas
enunciadas, alcanzadas en $x \equiv 0, \pi$. Media temporal en un
periodo $\tau = \tau(E_0)$: $x$ avanza exactamente $2\pi$, de modo que
\[
\frac1\tau\int_0^\tau y\,\dd t = \frac{x(\tau) -
x(0)}{\tau} = \frac{2\pi}\tau .
\]
La razón de velocidades extremas es $\sqrt{\frac{E_0 + 1}{E_0 -
1}} = 1 + O(E_0^{-1}) \to 1$: a energía alta, el rizo $\pm1$ del
potencial es despreciable frente a $E_0$, y el péndulo gira casi
uniformemente --- la tabla de lavar se aplana.

\textbf{25.} En $E_0 \geq 1 + \delta$, el integrando de
$\tau(E_0) = \int_{-\pi}^{\pi}\frac{\dd x}{\sqrt{2(E_0 +
\cos x)}}$ está dominado por $(2\delta)^{-1/2}$ y su derivada respecto
de $E_0$
\[
\partial_{E_0}\frac1{\sqrt{2(E_0 + \cos x)}}
= -\frac{1}{\bigl(2(E_0 + \cos x)\bigr)^{3/2}}
\]
por $(2\delta)^{-3/2}$, ambas \omterm{def:b3:lebesgue:l1}{integrables}
en el \omterm{def:b3:topology:compact}{compacto} $\intcc{-\pi}\pi$: la
derivación bajo la integral (el~\cref{thm:b3:lebesgue:paramdiff}) se
aplica y da $\tau'(E_0) < 0$ (el integrando es estrictamente negativo):
estrictamente decreciente, $\mathcal C^1$. Límites: cuando $E_0 \to
\infty$, $\tau \leq \frac{2\pi}{\sqrt{2(E_0 - 1)}} \to 0$; cuando
$E_0 \to 1^+$: escríbase $E_0 + \cos x = (E_0 - 1) +
2\cos^2\frac x2$; el integrando crece a \omterm{def:b3:measure:measure}{medida} que $E_0$ decrece, de
modo que, por convergencia monótona,
\[
\tau(E_0) \;\nearrow\;
\int_{-\pi}^{\pi}\frac{\dd x}{2\,\abs{\cos\frac x2}}
= +\infty,
\]
divergiendo la integral límite en $x = \pm\pi$ (allí
$\abs{\cos\frac x2} \sim \frac{\abs{x \mp \pi}}2$, a
no \omterm{def:b3:lebesgue:l1}{integrable} $\frac1{\abs\cdot}$): el
periodo explota al acercarse a la separatriz, en consonancia con el
$T(a) \to
\infty$ de la Parte II desde el lado de las libraciones. El eje de
energías se lee así: reposo en $E_0 = -1$; libraciones, $2\pi \nearrow
\infty$, en $\intoo{-1}1$; la separatriz infinitamente lenta en
$E_0 = 1$; rotaciones, $\infty \searrow 0$, más allá. Una integral, la
vida entera del péndulo.
\end{solution}
